\documentclass[a4paper,12pt]{article}
\usepackage[UTF8]{ctex} % 推荐使用 xelatex 编译
\usepackage{geometry}
\geometry{left=2.5cm,right=2.5cm,top=2.5cm,bottom=2.5cm}
\usepackage{graphicx}
\usepackage{caption}
\usepackage{booktabs}
\usepackage{tabularx}
\usepackage{float}
\usepackage{multirow}
\usepackage{array}
\usepackage{xcolor}

\begin{document}

% 第一页：封面
\begin{titlepage}
    \vspace*{1cm}
    \noindent 附件 1 \hfill \Large 序号：\underline{\makebox[3cm]{1674693}}\\[2.5cm]
    
    \begin{center}
        {\Huge \textbf{第十二届全国大学生光电设计竞赛}}\\[0.8cm]
        {\Huge \textbf{创意组}}\\[2.5cm]
        
        {\fontsize{36pt}{40pt}\selectfont \textbf{创意策划书}}\\[3cm]
        
        {\Large \textbf{竞赛题目：}\underline{\makebox[12cm]{动静有常--基于傅里叶频谱的脉搏实时采集分析仪}}}\\[1cm]
        {\Large \textbf{竞赛队名：}\underline{\makebox[12cm]{光脉守护者}}}\\[2cm]
    \end{center}
    
    \begin{table}[H]
        \centering
        \renewcommand{\arraystretch}{1.8}
        \begin{tabular}{|p{15cm}|}
        \hline
        \textbf{指导教师推荐理由：} \\
        \hline
        1、此项目将脉搏采集与傅里叶频谱分析相结合，集成了前端监测与后端处理分析，可作为家庭智能中医辅助诊断工具。 \\
        \hline
        2、此项目运用物联网技术，实现远程健康检测，将光电技术与新一代信息技术紧密结合。 \\
        \hline
        3、此项目利用了快速傅里叶变换（FFT）技术，能够高效、准确地处理和分析心电信号。 \\
        \hline
        4、此项目涉及Keil5软件编程、光电信息转换应用，符合参赛标准。 \\
        \hline
        \end{tabular}
    \end{table}
    
    \vfill
    \begin{center}
        {\Large \textbf{全国大学生光电设计竞赛委员会制}}\\[0.5cm]
        {\Large \textbf{二〇二四年四月}}
    \end{center}
\end{titlepage}

\newpage

% 成员表格（空表，仅做框架展示）
\begin{table}[H]
    \centering
    \renewcommand{\arraystretch}{1.5}
    \begin{tabular}{|c|c|c|c|c|c|c|}
        \hline
        \multicolumn{2}{|c|}{队长姓名} & & 学校 & & 学号 & \\
        \hline
        \multicolumn{2}{|c|}{专业班级} & \multicolumn{3}{c|}{} & 手机 & \\
        \hline
        \multicolumn{2}{|c|}{E-mail} & \multicolumn{5}{c|}{} \\
        \hline
        \multirow{5}{*}{合作者} & 姓名 & 学号 & 专业、班级 & \multicolumn{2}{c|}{所在学校} & 联系方式 \\
        \cline{2-7}
        & & & & \multicolumn{2}{c|}{} & \\
        \cline{2-7}
        & & & & \multicolumn{2}{c|}{} & \\
        \cline{2-7}
        & & & & \multicolumn{2}{c|}{} & \\
        \cline{2-7}
        & & & & \multicolumn{2}{c|}{} & \\
        \hline
        \multirow{2}{*}{指导教师} & 姓名 & 学校 & 所在院系 & \multicolumn{2}{c|}{联系电话/手机} & E-mail \\
        \cline{2-7}
        & & & & \multicolumn{2}{c|}{} & \\
        \hline
    \end{tabular}
\end{table}

\vspace{1cm}

\section*{一、产品的市场调研}
脉搏是一种周期性的生理性活动。脉搏信号包含丰富的生理信息，如心率、心律、血管状况等。脉搏信号的采集和分析在医疗诊断、健康管理、运动监测等领域具有重要意义。

根据市场调研，脉搏采集分析仪的市场规模和份额正在持续增长。2024年全球智能脉搏仪市场规模估计为13.4亿美元，预计到2029年将达到17.9亿美元，预测期内复合年增长率为6.15\%。

在中国，医用脉搏采集仪行业也在快速发展。2020年中国医用脉搏采集仪行业市场规模为10.8亿元，同比增长率为14.9\%，预计未来几年将继续保持较快的增长态势。这表明中国脉搏采集仪市场正在扩大，需求不断增长。此外，家用式脉搏采集仪在COVID-19疫情背景下需求显著增加，已成为家庭健康管理的重要工具之一。

随着远程医疗和自我健康管理理念的深入人心，开发一种将脉搏采集与傅里叶频谱分析相结合，集成了前端监测与后端处理分析的、完善的、家用健康监测器件具有重要的临床应用意义。

\section*{二、国内外技术分析}
脉搏采集技术在国内外的研究和发展呈现出多样化和技术进步的趋势。从传统的机械式测量到现代的数字化和自动化测量，技术不断进步，功能日益完善。这些技术的应用不仅提高了脉搏测量的精度，还为心血管健康监测和疾病诊断提供了更全面的信息。目前市场上的脉搏采集仪主要分为手指脉搏采集仪、手持式脉搏采集仪、手腕脉搏采集仪和远程脉搏采集仪等类型。智能脉搏采集仪市场的主要驱动因素包括心血管疾病患病率的上升、老年人口的增加以及对便携式监测设备的需求增长。这些设备主要用于医院、诊所、门诊手术中心和家庭医疗保健等领域。

与传统的家用脉搏测量仪器相比，本产品在对心电信号采集的同时，还能通过频谱特征得到实时健康数据，并能实现远程读取共享；而与医用级脉搏测量仪器相比，本产品成本低，可进行便携式测量。这使得我们的产品在前端监测和后端频谱特性分析方面具有优势，可作为家庭智能中医辅助诊断工具，能够为用户提供更准确的健康监测和分析结果，有望在市场上获得广泛的认可和应用。

\section*{三、成本核算}
预计花费 300 元左右。

\section*{四、技术方案论证}
我们的脉搏实时采集分析仪包括采集装置和上位机数据处理两部分，其中采集装置由心电采集模块，无线传输模块，物联网模块组成，能够实时采集出人体的心电信号。为了提高数据的准确性，我们首先在装置中植入了滤波算法对信号进行了初步的滤波处理，通过卡尔曼滤波器进行预测和更新两个步骤，不断调整状态估计，以提高对真实状态的估计精度。之后通过 keil 编程实现了心电信号（electrocardiogram, ECG）的采集。之后将处理后信号发送到上位机，在上位机中对心电信号进行快速傅里叶变换，得出变换后的频谱分布，据此频谱特性分析出人体运动状态的检测和判断。

\subsection*{(一) 采集装置}
\textbf{1. 控制核心：}以 STM32F103C8T6 为主控，控制各个模块的运行。\\
\textbf{2. 心电采集模块：}本设计心电测量芯片采用 ADS1292R 模拟前端芯片，ADS1292R 是一款高性能的模数转换器，主要用于将模拟信号转换为数字信号。ADS1292R 的输入端通过传输线连接三个电极片：RA、LA、RL。三个电极片分别接到右手腕、左手腕、右腹。由电极片采集的生物电压通过传输线经过芯片输出端传输到芯片内部，生物电压经过 ADS1292R 处理将模拟信号转换成数字信号，则此刻的数字信号作为心率信号。ADS1292R 的内部振荡器经过调整振荡信号按一定频率工作用于降噪，通过 SPI(串行外设接口)与 STM32 通信，通过滤波算法后得出初步的心电数据。

\textbf{(1) ADS1292R 电路设计}\\
图2为 ADS1292R 心电采集模块的电路设计原理图。

\begin{figure}[H]
    \centering
    \fbox{\begin{minipage}{0.9\textwidth}
    \centering
    \vspace{0.5cm}
    \textcolor{gray}{\LARGE [此处放置图片: ADS1292 心电采集模块电路图]}\\[0.5cm]
    \textbf{【图片内容】：} ADS1292芯片的外围引脚连接原理图，包含SPI通信接口及电极输入端设计。\\
    \textbf{【图片作用】：} 直观展示核心采集模块的底层硬件电路设计原理，论证项目在底层硬件电生理信号采集上的可行性与技术实现细节，体现项目硬件设计的专业度。\\
    \vspace{0.5cm}
    \end{minipage}}
    \caption{ADS1292 心电采集模块电路图}
\end{figure}

\textbf{(2) 心电信号采集方式}\\
将导联线的三个端子依次连接人体部位进行心电信号的采集，采集方式分为两种，一是黄色电极 L 贴在左胸口，红色电极 R 贴在右胸口，绿色电极贴在右侧小腹；二是黄色电极 L 贴在左手腕，红色电极 R 贴在右手腕，绿色电极贴在右侧小腹。我们采用第二种测量方式对被测者进行心电信号的采集，如图 2 所示。

\begin{figure}[H]
    \centering
    \fbox{\begin{minipage}{0.9\textwidth}
    \centering
    \vspace{0.5cm}
    \textcolor{gray}{\LARGE [此处放置图片: 导联线佩戴示意图及被测者测量示意图]}\\[0.5cm]
    \textbf{【图片内容】：} 左侧为人体标准导联电极贴放位置示意图，右侧为实际被测者手腕和腹部佩戴电极的照片。\\
    \textbf{【图片作用】：} 对文字描述的采集方式进行可视化补充，直观展示设备在实际测试和应用中的正确佩戴及可操作性。\\
    \vspace{0.5cm}
    \end{minipage}}
    \caption{导联线佩戴示意图及被测者测量示意图}
\end{figure}

\subsection*{(二) ECG 数据分析处理}
\textbf{1. 利用 MATLAB 进行 ECG 数据处理}\\
通过采集被测者心电信号，我们利用蓝牙无线将数据传输至上位机，并通过 MATLAB 进行编程，对采集到的被测者平静状态下的 ECG 信号进行分析处理。

首先将 CSV 文件导入数据，先对心电信号进行了归一化，得到了被测者的原始 ECG 心电信号图（图 4 所示）。之后，对其进行 FFT（Fast Fourier Transform，快速傅里叶变换）。其中，FFT 是一种用于将信号从时域转换到频域的数学算法。它是一种高效的数字信号处理技术，可以用来分析信号的频率成分。在信号处理中，一个信号通常被视为一系列在不同时间点上的值。这些值可以表示为时间的函数，即时域信号。FFT 的作用是将这样的时域信号转换为一个频域信号，其中频域信号是时域信号中不同频率成分的幅值和相位。具体来说，FFT 算法通过将时域信号分解为一系列正弦和余弦波，这些波具有不同的频率和振幅。这些正弦和余弦波的组合形成了原始信号。

通过将 FFT 的结果 Ayy 除以 N/2，以得到实际的幅度值。在 FFT 中，第一个元素对应于直流分量，而其余元素对应于不同频率的正弦波。由于直流分量通常需要单独处理，因此将 Ayy 的第一个元素除以 2，而其他元素则除以 N/2。然后将 Ayy 的第一个元素（直流分量）除以 2。由于直流分量通常不包含实际的频率信息，因此将其幅度减半。为避免在频谱分析中出现不必要的峰值，我们将 Ayy 的第一个元素（直流分量）的第一个值（幅度）设置为 0。

之后我们计算了实际的频率值 F。F 是一个向量，其元素对应于 Ayy 中每个元素的频率。F 的计算基于公式 Fn=(n-1)*Fs/N，其中 n 是 FFT 的索引，Fs 是采样频率，N 是 FFT 的点数。然后创建了一个幅度-频率曲线图，显示频率和对应的幅度的关系。

通过以上计算这些波的频率、振幅和相位，我们可以得到信号的频谱，即不同频率成分的贡献。此外我们计算了心率值，并有效识别和判断了异常的高频信号成分。如图 4 所示，我们分别显示了原始信号波形、心率值、相位幅度和频谱分布，以分析信号的频谱特性。

\begin{figure}[H]
    \centering
    \fbox{\begin{minipage}{0.9\textwidth}
    \centering
    \vspace{0.5cm}
    \textcolor{gray}{\LARGE [此处放置图片: matlab 程序分析代码截图]}\\[0.5cm]
    \textbf{【图片内容】：} 执行FFT变换、幅度转换、相位角度计算及绘制频谱图(`subplot`, `plot`)的核心MATLAB程序代码片段。\\
    \textbf{【图片作用】：} 展示核心数据分析算法的源码逻辑，佐证了项目利用FFT算法分析心电信号的技术真实性与可靠性。\\
    \vspace{0.5cm}
    \end{minipage}}
    \caption{matlab 程序中对原始信号进行分析的部分代码}
\end{figure}

\begin{figure}[H]
    \centering
    \fbox{\begin{minipage}{0.9\textwidth}
    \centering
    \vspace{0.5cm}
    \textcolor{gray}{\LARGE [此处放置图片: 信号的相位幅度和频谱分布图]}\\[0.5cm]
    \textbf{【图片内容】：} 三幅子图：脉搏实测信号时域波形（带BPM值）、FFT后的相位幅度散点图、FFT后的频谱分布曲线。\\
    \textbf{【图片作用】：} 直观展示算法处理前后的实际数据对比（时域波形至频域特征的转化），证明FFT算法在提取心跳特征上的有效性。\\
    \vspace{0.5cm}
    \end{minipage}}
    \caption{被测者的实测心电信号的相位幅度和频谱分布}
\end{figure}

此外，我们通过采集被测者运动前后的心电信号（图 5 所示），依据上述方法对脉搏和频谱进行分析比对，由图 5 可以发现：被测者运动前的脉搏处于相对稳定的状态，为每分钟 72 次，频谱基本分布低频区，基频信号为 0.6，不存在不正常的峰值或缺失以及异常的高频成分，而被测者运动后的心率增加，脉搏的频率增加，为每分钟 105 次，且出现了不规则的、短暂的波动，频谱基本分布仍然以低频区为主，基频信号为 0.875，且存在部分异常的高频成分（如运动后频谱中黑色圆圈所示）。

根据对被测者运动前后的频谱分析，我们可以评估被测者对该段运动的适应性和运动强度是否适宜。根据分析显示，被测者运动后频谱特征仍然保持稳定，没有明显的异常，这表明运动强度适宜，心脏能够有效应对。

相反，若被测者运动后的频谱中出现不正常的峰值或高频成分，这可能表明运动强度过大，心脏负担加重，需要调整运动强度或休息。

基于频谱分析的结果，通过监测运动前后的脉搏和频谱分布的变化，我们可以评估心血管系统对运动的反应，从而制定更有效的预防和治疗策略。我们可以制定个性化的运动训练计划，确保运动强度适宜，避免过度训练，从而达到最佳的心血管健康促进效果。此外，频谱分析还可以用于长期的健康监测，通过定期记录和分析心电信号，我们可以监测心血管健康状况的变化，及时发现潜在的健康问题，对于预防心血管疾病具有重要的意义。

\begin{figure}[H]
    \centering
    \fbox{\begin{minipage}{0.9\textwidth}
    \centering
    \vspace{0.5cm}
    \textcolor{gray}{\LARGE [此处放置图片: 运动前后脉搏和频谱分布对比图]}\\[0.5cm]
    \textbf{【图片内容】：} 包含四幅子图，上下并排横向对比了被测者“运动前”与“运动后”的心电时域波形和频谱分布变化，并用红圈标出了异常高频成分。\\
    \textbf{【图片作用】：} 利用真实采集数据直接证明产品的核心功能——“评估运动对心电信号影响及健康状态判断”的可行性与科学依据。\\
    \vspace{0.5cm}
    \end{minipage}}
    \caption{被测者运动前后的脉搏和频谱分布图}
\end{figure}

\textbf{2. 利用物联网模块使处理结果可视化}\\
物联网模块中采用 ESP8266-01S 芯片作为 WIFI 模块，STM32 通过串口向 ESP8266 发送 AT 指令接受到 AT 指令后，利用 MQTT 协议连接到阿里云物联网平台，ESP8266 作为发布者，向阿里云发送数据，阿里云作为订阅者，接受 ESP8266 发送来的数据，利用阿里云平台的 IOT Studio 服务可视化开发绘制 ECG 心电数据显示界面（图 7 所示）。

\begin{figure}[H]
    \centering
    \fbox{\begin{minipage}{0.9\textwidth}
    \centering
    \vspace{0.5cm}
    \textcolor{gray}{\LARGE [此处放置图片: MQTT连接程序的代码截图]}\\[0.5cm]
    \textbf{【图片内容】：} STM32通过ESP8266配置MQTT主题订阅（Subscribe）和发布（SendMessage）的具体C语言功能代码。\\
    \textbf{【图片作用】：} 展示物联网模块连接云端的核心代码逻辑，证明设备具备数据上云与远程传输的能力，是产品“物联网技术”亮点的重要支撑。\\
    \vspace{0.5cm}
    \end{minipage}}
    \caption{关于 esp8266 通过 MQTT 连接阿里云的程序}
\end{figure}

\begin{figure}[H]
    \centering
    \fbox{\begin{minipage}{0.9\textwidth}
    \centering
    \vspace{0.5cm}
    \textcolor{gray}{\LARGE [此处放置图片: 阿里云 IOT Studio 可视化界面图]}\\[0.5cm]
    \textbf{【图片内容】：} 网页端的UI界面，上方显示实时的心电电压数值，下方显示持续更新的ECG模拟电压值和转换电压值的双线折线图。\\
    \textbf{【图片作用】：} 展示最终的数据上云效果和前端UI可视化结果，体现项目的完整技术闭环以及向终端用户提供直观数据呈现的能力。\\
    \vspace{0.5cm}
    \end{minipage}}
    \caption{利用阿里云 IOT Studio 绘制的简单显示界面}
\end{figure}

\section*{五、技术路线设计}
本设计主要研究内容为设计一个脉搏实时采集分析仪，用于帮助人们评估心血管系统对运动的反应，监测心血管健康状况。
\begin{enumerate}
    \item 设计心电采集电路，集成了心电采集所需部件，方便设备小型化。
    \item 运用 Keil5 编程，使用低成本单片机控制各个模块实现对应功能。
    \item 获取正常人体运动前后心电数据，分析其特征。
    \item 利用 MATLAB 编程，将待测者心电信号进行频谱分析，将分析结果与正常人结果对比，判断待测者的健康状况。
    \item 使用物联网芯片将产品连接到云平台，实现远程监测与数据共享，增加产品实用性。
\end{enumerate}

\section*{六、产品的外观设计}
基于傅里叶频谱的脉搏实时采集分析仪主要包括采集电路模块、数据分析处理模块和无线传输模块，物联网模块。

\begin{figure}[H]
    \centering
    \fbox{\begin{minipage}{0.9\textwidth}
    \centering
    \vspace{0.5cm}
    \textcolor{gray}{\LARGE [此处放置图片: 产品实物图]}\\[0.5cm]
    \textbf{【图片内容】：} 组装完成的硬件实物照片，能看到绿色的底层PCB连接着STM32核心板、ADS1292模块、ESP8266模块及OLED显示屏等组件。\\
    \textbf{【图片作用】：} 展示项目的硬件实物拼装最终形态，证明创意已具备实物原型雏形，大大提高了整个策划书的可信度和项目落地的可行性。\\
    \vspace{0.5cm}
    \end{minipage}}
    \caption{产品实物图}
\end{figure}

\section*{七、产品的营销策划}
\textbf{1. 市场定位与目标客户：}\\
① 定位：我们的产品是一款集成了 STM32F103CT86 微控制器、ADS1292R 模数转换器、OLED 显示屏、ESP8266 模块、蓝牙 BLE 模块和电脑端软件的集成化健康监测与分析平台。\\
② 目标客户：\\
1) 个人健康监测用户：本产品可进行长期的健康监测，通过定期记录和分析心电信号，我们可以监测心血管健康状况的变化，及时发现潜在的健康问题。\\
2) 运动爱好者：本产品可监测运动前后的脉搏和频谱分布的变化，我们可以评估心血管系统对运动的反应，制定个性化的运动训练计划，确保运动强度适宜，避免过度训练，从而达到最佳的心血管健康促进效果。\\
3) 老年人等不便行动者：随着年龄的增长，老年人的心血管系统可能会出现各种问题。本产品可以作为家庭智能中医辅助诊断工具，帮助家庭成员和护理人员监测老年人的健康状况，及时发现潜在的健康风险。

\textbf{2. 产品优势与创新点：}\\
① \textbf{该产品集成性高，构建了一个完整的心电信号采集和傅里叶频谱分析系统。}相比传统家用脉搏监测仪，该作品通过将脉搏采集与傅里叶频谱分析相结合，形成了一个集成了前端监测与后端处理分析的、完善的、家用健康监测器件，可靠性高。\\
② \textbf{采用 FFT 算法对心电信号进行频谱分析}，相较于传统的离散傅里叶变换（DFT），FFT 的计算效率更高，能够更快速地处理大量数据。通过对运动前后的心电数据进行对比分析，该作品能够评估运动对心电信号的影响，为运动健康监测和个性化训练提供科学依据。\\
③ \textbf{无线数据链通信传输，支持多设备连接通信。}设备灵活性高，轻巧便携，能够将设备与手机端，电脑端无缝连接，同时数据可导出，实现多平台数据共享，能够适应不同的应用场景和通信需求。\\
④ \textbf{运用物联网技术，将数据上传云端并记录保存，可实现远程医疗诊断。}心电数据的云端存储和远程访问，提高了数据的安全性和可访问性，为远程医疗和健康管理提供了便利。

\textbf{3. 营销策略：}\\
① 品牌宣传：通过线上和线下渠道，加强品牌宣传，提高品牌知名度和影响力。\\
② 产品推广：在各大电商平台、医疗器械展会、健康论坛等平台进行产品推广，突出产品优势和创新点。\\
③ 合作伙伴：与医疗机构、远程医疗平台、运动品牌等建立合作伙伴关系，共同推广产品。\\
④ 用户体验：通过线上线下活动，邀请目标客户体验产品，收集用户反馈，不断优化产品性能。\\
⑤ 售后服务：提供专业的技术支持和售后服务，确保用户在使用过程中无后顾之忧。

\textbf{4. 营销渠道：}\\
① 线上渠道：电商平台、官方网站、社交媒体、短视频平台等。\\
② 线下渠道：医疗器械展会、健康论坛、医疗机构合作等。

\section*{八、团队情况}
本团队成员有指导教师一名，在校本科生两名，在校研究生一名。

项目负责人主要负责：项目整体策划与管理，协调团队成员的工作分配和进度；并实施傅里叶频谱分析的算法编程，对处理的 ECG 心电信号进行全面评估。以及撰写项目报告和申报材料。

成员 1 主要负责：设计产品的心电采集电路，搭建好了硬件环境。编写了 STM32 控制器程序和物联网模块的程序，为芯片植入了中值滤波，卡尔曼滤波初步处理算法。

成员 2 主要负责：进行市场调研，分析市场规模和潜在客户群体并制定产品的营销策略和推广计划。对产品进行测试，收集数据，并根据测试结果进行产品优化。

\end{document}
